\section{结果}\label{sec:results}

本节按固定顺序报告证据：窗口、它的时序、机制、迁移、边界，以及现代译码器参照。

\subsection{首轮是一个双向窗口}\label{sec:formal}

每个实例由 C0 与只在首轮阈值加 $+2$ 的 FRO-2 分别译码。两个变体看到同一把密钥、同一个错误与同一个综合征。对比在 $P_1'=(4801,45,92)$ 上进行，覆盖 RAND、TYPE-I 与 NCW-shallow。表~\ref{tab:main} 给出计数与检验；图~\ref{fig:main} 给出绝对率、区间与逐钥分布。

FRO-2 在三种条件下都降低失败率（表~\ref{tab:main}）。RR 为 4.15 [3.29, 5.38]、5.47 [4.25, 7.28] 与 4.55 [3.21, 6.82]。每个区间都排除 1，密钥聚类符号检验在 Holm 校正后也一致。逐钥看，FRO-2 在 278 个密钥-条件情形中改善 233 个、变差 18 个、持平 27 个。预指定门槛要求最差条件的 $95\%$ DFR 上界至少改善 $0.5$ 个 $\log_{10}$ 单位（约 3.2 倍）。全池改善 0.679 个 $\log_{10}$（约 4.8 倍），纯新子集改善 0.687 个 $\log_{10}$，因此门槛通过。

窗口两侧并不对等。FRONEG 在首轮减 1：它使 92 把 RAND 密钥全部变差，失败数从 282 升到 2{,}235，RR 为 0.126 [0.111, 0.141]，即上升 7.93 倍。向下走一步的代价因此超过向上走两步的收益。

完整局部扫描分辨出两侧（图~\ref{fig:window}）。偏移 $+1$ 至 $+3$ 使失败率下降约 2.3 至 5.5 倍，最大值落在 $+2$ 或附近。宽松侧在两个计数内崩塌：降低 1 给出 RR 0.10--0.13，降低 2 给出 0.0014--0.0030。在 RAND 下，两计数下调把失败数放大 365 倍，从 276{,}000 次中的 282 次到 102{,}980 次。窗口是双向且陡峭不对称的。直白地说，保守偏移换来 4--5 倍的下降，而向下走一步的代价超过这一收益。

\begin{table*}[!t]
\centering
\small
\caption{$P_1'=(4801,45,92)$ 处的主对比。单元格给出失败数/译码数（DFR）。RR 为 C0 除以 FRO-2，带 95\% 密钥聚类自助区间；Holm $p$ 为跨三条件的密钥聚类符号检验。}
\label{tab:main}
\begin{tabular}{l r r r r r}
\toprule
条件 & 密钥数 & C0 失败数/译码数 (DFR) & FRO-2 失败数/译码数 (DFR) & RR [95\% CI] & Holm $p$ \\
\midrule
RAND & 92 & 282/276{,}000 ($1.02\times10^{-3}$) & 68/276{,}000 ($2.46\times10^{-4}$) & 4.15 [3.29, 5.38] & $4.88\times10^{-17}$ \\
TYPE-I & 94 & 372/282{,}000 ($1.32\times10^{-3}$) & 68/282{,}000 ($2.41\times10^{-4}$) & 5.47 [4.25, 7.28] & $3.05\times10^{-22}$ \\
NCW-shallow & 92 & 141/276{,}000 ($5.11\times10^{-4}$) & 31/276{,}000 ($1.12\times10^{-4}$) & 4.55 [3.21, 6.82] & $7.66\times10^{-12}$ \\
\bottomrule
\end{tabular}
\end{table*}

\dblfig{fig2_main_results.pdf}{$P_1'=(4801,45,92)$ 处的主对比。(a)~每 10{,}000 次译码的绝对失败数。(b)~C0/FRO-2 风险比与 95\% 密钥聚类自助区间，虚线为 1。(c--e)~RAND、TYPE-I 与 NCW-shallow 下逐钥避免失败数的分布；标注给出改善、变差与持平的密钥数。每把密钥每条调度 3{,}000 次译码。}{fig:main}

\begin{figure*}[!t]
\centering
\includegraphics[width=0.88\textwidth]{fig6_threshold_window.pdf}
\caption{$P_1'$ 处首轮是一扇双向窗口。(a)~六个偏移的绝对 DFR（对数轴）。(b--d)~正偏移 $+1$、$+2$、$+3$ 按条件给出的风险比，统一从 0 起的线性轴，带 95\% 密钥聚类区间；浅色带只标出 $+2$，不表示统计优越性。负偏移的绝对率见 (a)。}
\label{fig:window}
\end{figure*}

\subsection{偏移的轮次与持续方式}\label{sec:temporal}

轮次对照改变偏移的施加轮次与持续方式。R2one 只在第 2 轮加 $+2$；R2plus 从第 2 轮起持续加 $+2$。同样的幅度在首轮给出 RR 4.15，在第 2 轮给出 1.21，持续到后续轮则给出 0.0168。

在 RAND 下，R2one 给出 1.21 [1.07, 1.36]，44 把密钥改善、25 把变差、23 把持平。R2plus 给出 0.0168 [0.0149, 0.0187]：失败数从 282 升到 16{,}785，DFR 升到 $6.08\times10^{-2}$，92 把密钥中 0 把改善。TYPE-I 与 NCW-shallow 重复同一模式：R2one 给出 1.45 与 1.34，R2plus 给出 0.0210 与 0.0118，94 把与 92 把中均 0 把改善。首轮减一对照保持有害符号，为 0.109 与 0.101。

扩大确认在 P2 与 P3 重复同一检验（表~\ref{tab:temporal}，图~\ref{fig:temporal}）。六个条件条目中，首轮偏移的 RR 跨 1.76--4.62（FRO-2 为 2.37--4.62），第 2 轮一次性偏移跨 1.13--1.54，持续偏移跨 0.009--0.019。没有哪个偏移处处最优：FRO-1 领跑 P2 的 RAND 列，FRO-2 领跑其余五列。第 2 轮区间在六个单元中有四个单元排除 1。轮次对照变体在主分析之后加入，属于探索性结果；扩大确认的调度是预指定的。决定是否有益的是偏移的轮次与持续方式，而不是单独的幅度。

\begin{table*}[!t]
\centering
\small
\caption{P2 $(3001,37,66)$ 与 P3 $(4001,41,82)$ 上的扩大时序确认。C0 列为绝对 DFR；其余为风险比（C0 除以该调度），区间见图~\ref{fig:temporal}。每格 50 把密钥、每把 1{,}000 个错误。}
\label{tab:temporal}
\begin{tabular}{l r rrrrr}
\toprule
测试单元 & C0 DFR & FRO-1 & FRO-2 & FRO-3 & R2-2 & UNIFORM-2 \\
\midrule
P2, RAND & $2.56\times10^{-3}$ & 2.56 & 2.37 & 1.91 & 1.13 & 0.0135 \\
P2, TYPE-I & $3.14\times10^{-3}$ & 2.31 & 2.85 & 2.15 & 1.29 & 0.0188 \\
P2, NCW-shallow & $1.20\times10^{-3}$ & 1.76 & 3.00 & 2.50 & 1.36 & 0.00949 \\
P3, RAND & $1.10\times10^{-3}$ & 2.29 & 2.62 & 2.62 & 1.34 & 0.00951 \\
P3, TYPE-I & $1.70\times10^{-3}$ & 3.27 & 3.27 & 2.18 & 1.31 & 0.0165 \\
P3, NCW-shallow & $7.40\times10^{-4}$ & 3.70 & 4.62 & 3.08 & 1.54 & 0.00925 \\
\bottomrule
\end{tabular}
\end{table*}

\dblfig{fig4_temporal_dissociation.pdf}{时序解离。(a--c)~RAND、TYPE-I 与 NCW-shallow。上排：$P_1'$ 的主点轮次对照（灰色菱形）。下排：迁移点 P2（实心圆）与 P3（空心方块）。区间为 95\% 归档区间；虚线为风险比 1。}{fig:temporal}
